\section{The Thorin approach in summary, and a path to RH}\label{ch:path}

The path is a programme with its gap identified, not a proof.

\subsection{The path along the basepoint}\label{sec:ch25:path}

Figure~\ref{fig:ch25:path} and Table~\ref{tab:ch25:path} set out the path.

\begin{figure}[htbp]
\centering
\begin{tikzpicture}[node distance=0.55cm and 0.5cm, every node/.style={font=\small}, box/.style={draw, rounded corners, align=center, minimum width=3.4cm, minimum height=0.9cm, inner sep=3pt}, arr/.style={->, thick}]
\node[box, fill=blue!8] (primes) {primes, gamma factor,\\ pole (L\'evy side)};
\node[box, fill=orange!12, right=1.3cm of primes] (lemma) {sine squared lemma\\ (positive kernel)};
\node[box, fill=blue!8, right=1.3cm of lemma] (zeros) {zeros as atoms\\ (Thorin side)};
\node[box, fill=green!10, below=of lemma] (rho) {Thorin density $\rho_\varepsilon$\\ at basepoint $\frac12+\varepsilon$};
\node[box, fill=yellow!15, below=0.8cm of rho] (amid) {$\frac12<\alpha<1$: GGC iff\\ no zeros right of $\alpha$;\\ almost GGC};
\node[box, fill=green!10, left=0.8cm of amid] (a1) {$\alpha\ge1$: GGC\\ proved (Chapters~\ref{chap:basepoint} to \ref{chap:almost})};
\node[box, fill=red!10, right=0.8cm of amid] (acentre) {$\alpha\downarrow\frac12$: GGC\\ iff RH};
\node[box, fill=gray!10, below=0.7cm of amid] (need) {needed: $T_\alpha$ a GGC at one $\alpha<\frac78$,\\ then uniformly in $\alpha$ and height};
\draw[arr] (primes) -- (lemma); \draw[arr] (lemma) -- (zeros);
\draw[arr] (lemma) -- (rho); \draw[arr] (zeros) |- (rho);
\draw[arr] (rho) -- (a1); \draw[arr] (rho) -- (amid); \draw[arr] (rho) -- (acentre);
\draw[arr] (a1) -- (amid); \draw[arr] (amid) -- (acentre);
\draw[arr] (amid) -- (need);
\end{tikzpicture}
\caption{The Thorin path to RH. The sine squared lemma connects the arithmetic to the Thorin density. The programme must carry the GGC property across the strip.}
\label{fig:ch25:path}
\end{figure}

\begin{table}[htbp]
\centering\small
\begin{tabular}{p{0.05\textwidth}p{0.55\textwidth}p{0.3\textwidth}}
\toprule
step & statement & status\\
\midrule
1 & The clock is infinitely divisible and self decomposable; the heat trace is positive and decreasing & proved (Theorem~\ref{thm:ch12:idsd}, which is Theorem A of \cite{P26}), with the verification of RH below height $1000$ as input\\
2 & The heat trace is log convex, its leading Hankel blocks of size up to $3\times10^7$ are positive definite, derivative signs hold to order $1.16\times10^{25}$ & proved, with verification height as input\\
3 & The reciprocal at basepoint $\alpha\ge1$ is a GGC with Poisson Thorin density, and at $\alpha=1$ its L\'evy form, hitting time, scale mixture and completeness are explicit; for every $\alpha>\frac12$ the shifted clock $T_\alpha$ exists and is an almost GGC & proved (Theorems~\ref{thm:poisson}, \ref{bp:thm:alpha} and \ref{ag:thm:almost}; Chapters~\ref{chap:basepoint} to \ref{chap:almost}), with the verification of RH below height $1000$ as input to Theorem~\ref{bp:thm:alpha}\\
4 & The normalised reciprocal zeta function is infinitely divisible; the canonical function of the floor of the reciprocal Xi function is an arithmetic part plus a completely monotone archimedean part, and RH is the complete monotonicity of the sum & proved (Chapters~\ref{chap:rzeta} and \ref{chap:rzid})\\
5 & The Thorin density is computable at every resolution, with an explicit bound for each source of error, and is positive at every sample point of every window computed; the runs do not yet add these bounds to the ball radii (Section~\ref{st:sec:wp2}) & computed (Chapters~\ref{chap:evaluation} to \ref{chap:heat})\\
6 & Zero density $\nu(\frac12+\varepsilon,200)=0$ for every computed $\varepsilon$, down to $\varepsilon=0.005$ & computed (Section~\ref{ch:turing})\\
7 & For some $\alpha<1$ the shifted clock $T_\alpha$ is a GGC, equivalently, the half plane $\re s>\alpha$ is zero free (Theorem~\ref{thm:poisson}); then for every $\alpha'>\alpha$ the measure $U_{\alpha'}$ is the Thorin measure of $T_{\alpha'}$ and is positive at all heights & announced for $\alpha=\frac78$ in October 2026, with a formal proof in Lean, not yet refereed (Section~\ref{sec:ch25:open}); open for every $\alpha<\frac78$\\
8 & Positivity of $U_\alpha$ for every $\alpha>\frac12$ (Lagarias), equivalently the Thorin condition & open; equivalent to RH\\
\bottomrule
\end{tabular}
\caption{The path, step by step.}
\label{tab:ch25:path}
\end{table}



Section~\ref{bp:sec:compare} draws the two conclusions that fix the shape of the path. At the basepoint one the constructions attached to the reciprocal Xi function are unconditional and explicit (Sections~\ref{ch:basepoint}, \ref{ch:bphit} and \ref{ch:strip}, Table~\ref{sb:tab:results}), with two exceptions: the Krein string representation of Corollary~\ref{bp:cor:kent} assumes RH and simple zeros, and Lemma~\ref{bp:lem:wald} gives the Wald defect $\sum_\rho\rho^{-1}/\sqrt2$ only for a boundary bounded at infinity, while the boundedness of the computed boundary is observed, not proved. Below the basepoint one each construction holds, up to a boundary case, exactly when no zero lies to the right of $\alpha$ (Theorems~\ref{bp:thm:alpha} and \ref{sm:thm:strip}, Propositions~\ref{bp:prop:below} and \ref{wr:prop:wiener}), and the theory of Sato and Bondesson reaches nothing there (Section~\ref{sb:sec:cannot}).

\subsection{The open steps}\label{sec:ch25:open}

By Theorem~\ref{thm:poisson}, that is \cite[Theorem 3.7 and Remark 3.10]{P26}, step 7 is exactly a zero free half plane $\re s>\alpha$ with $\alpha<1$, a quasi Riemann hypothesis. Until 2026 no such half plane was known. In October 2026 OpenAI released a proof, produced with one of its AI models, that $\zeta(s)\ne0$ for $\re s>\frac78$, and the same for every Dirichlet $L$ function \cite{OAI26a}, together with a formal proof of both statements in the Lean proof assistant and a separate proof of the weaker half plane $\re s>\frac{11}{12}$ \cite{OAI26b}. The argument proves a power saving for smoothed sums of the M\"obius function twisted by Hecke characters of $\mathbb Q(\sqrt{-3})$, and it uses nothing from the Thorin picture. It has not yet been refereed. We rebuilt the formal proof from the public repository (commit \texttt{adc7f12}, Lean 4.34.1, Mathlib \texttt{d13f23b}): the Lean kernel accepts both proofs, which use only the axioms \texttt{propext}, \texttt{Classical.choice} and \texttt{Quot.sound}, and the comparator of the Lean Focused Research Organization, run without its sandbox, confirms that the theorems proved are the published challenge statements. If it stands, step 7 holds with $\alpha=\frac78$, and the constructions of this book follow. By Theorem~\ref{thm:poisson} the shifted clock $T_\alpha$ is a GGC for every $\alpha\ge\frac78$, with Thorin density \eqref{eq:Ualpha} and $v_\alpha>0$ for $\alpha>\frac78$, while at $\alpha=\frac78$ each zero $\frac78+i\gamma$ with $\gamma>0$, which the announcement does not exclude, would add an atom at $\gamma^2$ (Theorem~\ref{rz:thm:threshold}); for $\alpha>\frac78$ the clock $T_\alpha$ is a Brownian first passage time plus an independent GGC with no drift, because the first two steps of the proof of Proposition~\ref{bp:prop:valpha} go through for $\frac78<\alpha<1$, where every zero has $0<\alpha-\beta<\alpha$, with smallest margins $0.0137$ and $0.0157$ in ball arithmetic (\texttt{valpha\_min.py}); by Proposition~\ref{bp:prop:below} the compensated L\'evy form converges for $\frac78<\alpha<1$, with error $O(X^{7/8-\alpha+\delta})$ at prime cutoff $X$ for every $\delta>0$, and the same holds for $L(s,\chi_{12})$ (Section~\ref{bp:eta:sec:below}); by Theorem~\ref{ri:thm:ggcriesz} Riesz's function satisfies $R(x)=O(x^{7/16+\delta})$ for every $\delta>0$; and by de Bruijn's theorem $\Lambda\le\frac9{32}$ (Section~\ref{ch:dbn}), weaker than the bound $0.2$, while for the cosine tilts of Section~\ref{fm:sec:tilts} his theorem for cosine factors gives $\Omega\le\frac38$, against $\frac12$ with no hypothesis. Below $\frac78$ nothing changes. Positivity of the Thorin density on the single line $\re s=\alpha$ is a weaker statement than a zero free half plane (Remark~\ref{rem:ch2:oneline}). So the first theorem the Thorin path needs beyond the announcement, a zero free half plane $\re s>\alpha$ with $\alpha<\frac78$, is an open problem in its own right; the L\'evy form, the scale mixture, completeness and Riesz's criterion (Sections~\ref{ch:basepoint}, \ref{ch:strip} and \ref{ch:riesz}) restate it and do not get round it. The Thorin formulation offers two specific tools for it. The sine squared form, compensated as in Proposition~\ref{bp:prop:below}, writes the required positivity as domination of the pole piece by the prime and archimedean pieces at basepoint $\alpha$, a statement about a single explicit function of $\theta$ for each $\alpha$; and the Poisson form shows that a zero to the right of $\alpha$ appears as a Cauchy kernel of negative width, whose signature Section~\ref{ch:offline} calibrates exactly.

Step 8 carries the GGC property from any fixed $\alpha$ to every $\alpha>\frac12$. Weak limits cannot do this on their own (Section~\ref{ch:weak}), because the limit of the Cauchy flow is positive whether or not RH holds. Three mechanisms could do it, each consistent with everything computed in this book:
\begin{enumerate}
\item a positive approximating family whose positivity is proved independently, such as zeta functions of Beurling generalised prime systems that provably have only real ordinates, finite Krein strings, or the partial products over verified zeros completed by the Weyl tail; once the first $N$ zeros are verified, RH is equivalent to the GGC property of the tail factor left when their factors are removed from the clock \cite[Corollary 18.11]{P26};
\item an operator whose spectrum is the set of ordinates, where positivity would follow from self adjointness, the Hilbert P\'olya idea in Thorin form (Section~\ref{sec:ch18:kent});
\item a direct inequality for the variogram pairing that controls the pole by the primes and the gamma factor at every scale, the Weil form of step 8.
\end{enumerate}

\subsection{What computation contributes}\label{sec:ch25:contrib}

Computation cannot supply the uniformity in the height that steps 7 and 8 require: neither finitely many heat trace tests (Theorem~\ref{thm:ch13:perturb}) nor positivity at a fixed resolution (Section~\ref{sec:ch13:visibility}) can exclude an off line zero above the verified height. What it can do is developed in Section~\ref{sec:open}. It calibrates the size of the positivity margins, which tells a proof how much room it has: the Thorin density at resolution $\varepsilon$ has minimum $0.0147\,\varepsilon$ on the window below height $200$ (Section~\ref{ch:motivation}), a margin that vanishes linearly and must be controlled uniformly. That this margin is the value at the origin is a theorem on the whole verified window, and at every height for every $\varepsilon<\frac12$ it is equivalent to RH (Theorem~\ref{thm:dbeta:minall}).

Section~\ref{ch:sin2more} gives the variogram form of the domination inequality, on which the linear programming formulation of \cite[Sections 18.2.8 and 18.2.9]{P26} works, the most concrete target for step 8. Both are posed at the centre, where the variogram form is Weil's criterion and the linear program asks whether the clock $1/G$ is a GGC, so each is equivalent to RH. Section~\ref{ch:weak} and \cite[Section 18.2.3]{P26} give the closure criterion in its weakest form, the target for approximating families in step 8. Section~\ref{ch:almost} closes with four questions between almost and exact positivity (Section~\ref{ag:sec:exactly}), the last of which is step 7 below $\frac78$, the GGC property at a single basepoint $\alpha<\frac78$.
